\section{Scope, Method, and Qualifications}

\subsection*{Subject and Intended Reader}

This study is Part I of a four-part Catholic series on the Four Last Things (Death, Judgment, Hell, and Heaven). It is written for serious Catholic believers seeking a sober, source-grounded exposition of the Church's doctrine and spiritual tradition on death, moving from Sacred Scripture, conciliar definitions, Fathers, and Scholastic Doctors through the martyrs, the \textit{Ars Moriendi}, the last sacraments, and the burial liturgy into nightly Compline, examination of conscience, and received vocal prayer. Questions belonging primarily to Particular and General Judgment, Purgatory, and Indulgences (Part II, \textit{The Books Were Opened}), Hell (Part III, \textit{Depart from Me}), or the Beatific Vision and the glorified body (Part IV, \textit{Come, Ye Blessed}) are touched only insofar as they directly bound the doctrine of bodily death.

\subsection*{Scripture, Liturgical Texts, and Received Prayers}

All English quotations of Sacred Scripture follow the Challoner revision of the Douay-Rheims Bible, paired where helpful with the Clementine Vulgate Latin. Verse numbering follows the Vulgate and Douay-Rheims convention throughout, with divergences from Hebrew or modern critical numbering noted where they could otherwise confuse a reader---chiefly Psalm 115:15 in Douay-Rheims (Ps.\ 115:6 in the Clementine Vulgate; Ps.\ 116:15 in the Hebrew Psalter) and 1 Thessalonians 4:12--17 in Douay-Rheims (1 Thess.\ 4:13--18 in the Vulgate and modern editions).

Latin liturgical texts---the Ash Wednesday imposition of ashes (\textit{Memento, homo}) and Introit (\textit{Misereris omnium}), and the Easter Sequence (\textit{Victimae paschali laudes})---follow the 1962 \textit{Missale Romanum}; the Requiem Preface for the Dead (\textit{Praefatio pro Defunctis}) is verified directly in the 1922 Tours (Mame) edition of the \textit{Missale Romanum}, whose Latin wording agrees with the 1962 typical edition. The Office of the Dead (\textit{Placebo}, \textit{Dirige}, and the Responsory \textit{Credo quod Redemptor meus vivit}) is verified in the 1915 Benziger bilingual edition. The classical pre-1972 \textit{Rituale Romanum} rites of Holy Viaticum, Extreme Unction, Recommendation of a Departing Soul (\textit{Proficiscere}, \textit{Commendo te}, \textit{Libera, Domine}, \textit{Subvenite}), and Burial of the Dead (\textit{Non intres}, \textit{In paradisum}, \textit{Ego sum}), together with the pre-1911 Roman Breviary Evening Compline, are verified directly in the 1889 \textit{Manual of Prayers for the Use of the Catholic Laity}, published by order of the Third Plenary Council of Baltimore. Where liturgical discipline changed across reforms, those historical boundaries are stated beside the rites: the pre-1972 / 1962 \textit{Rituale Romanum} discipline alongside Pope St.\ Paul VI's Apostolic Constitution \textit{Sacram unctionem infirmorum} (1972) and the postconciliar \textit{Pastoral Care of the Sick}; and the pre-1911 four-psalm Compline alongside Pope St.\ Pius X's 1911 Psalter reform (\textit{Divino afflatu}) in the 1962 Roman Breviary and Pope St.\ Paul VI's 1970 \textit{Laudis canticum} promulgating the 1971 \textit{Liturgia Horarum}.

Every text presented for vocal prayer in Section 10 reproduces a received public-domain edition verbatim: the nightly Compline texts (Short Responsory, Canticle of Simeon with antiphon, and Collect, omitting the intervening \textit{Preces}) and the \textit{Anima Christi} from the 1889 \textit{Manual of Prayers} (pp.\ 73--75, 103), and St.\ Alphonsus Maria de Liguori's prayer in preparation for death from the 1843 Dublin Duffy edition of \textit{Preparation for Death} (Fourth Consideration, Third Point, ``Affections and Prayers,'' pp.\ 68--69), resolving only three obvious nineteenth-century OCR printing errors (\textit{Jast} to \textit{last}, \textit{hfe} to \textit{life}, and terminal \textit{P} to \textit{?}) without altering wording. Brief translations of Latin liturgical or patristic phrases embedded inside analytical prose serve solely as explanatory aids.

\subsection*{Source Hierarchy, Editions, and Limits}

The exposition distinguishes degrees of doctrinal authority, consolidated in Appendix A: Sacred Scripture and solemn dogmatic definitions (\textit{de fide}); ordinary magisterial teaching; patristic and scholastic consensus (\textit{sententia communis}); and open theological questions on which approved Catholic Doctors differ (such as the precise mode of natural cognition in the separated soul before the Resurrection).

Patristic, scholastic, medieval, and spiritual works are quoted directly from identified public-domain editions: Roberts--Donaldson--Coxe's \textit{Ante-Nicene Fathers} (St.\ Ignatius of Antioch, shorter recension; St.\ Cyprian, \textit{De mortalitate}), Philip Schaff and Henry Wace's \textit{Nicene and Post-Nicene Fathers} (St.\ Athanasius, \textit{On the Incarnation}; St.\ Cyril of Jerusalem, \textit{Catechetical Lectures}; St.\ Ambrose, \textit{De excessu fratris Satyri}; St.\ John Chrysostom, \textit{Homilies on Philippians} and \textit{First Thessalonians}; St.\ Augustine, \textit{Confessions}; St.\ Leo the Great, \textit{Sermons}; St.\ John of Damascus, \textit{De fide orthodoxa}), Marcus Dods's 1871 translation of \textit{The City of God} (collated with E.\ Hoffmann's CSEL 40 Latin), the English Dominican translation of St.\ Thomas Aquinas's \textit{Summa Theologiae} (Prima, Prima Secundae, and Tertia Pars), Frances M.\ M.\ Comper's 1917 edition of \textit{The Book of the Craft of Dying}, John O'Sullivan's translation of St.\ Robert Bellarmine's \textit{Commentary on the Book of Psalms}, J.\ Donovan's 1829 translation of the \textit{Catechism of the Council of Trent}, the Frederick Pustet Catholic edition of St.\ Francis de Sales's \textit{Introduction to a Devout Life}, and the 1843 Duffy \textit{Preparation for Death}. Conciliar and magisterial texts held as cataloged bibliographic references---the Councils of Vienne (DS 902), Florence (\textit{Laetentur caeli}, DS 1304--1306), Fifth Lateran (\textit{Apostolici Regiminis}, DS 1440), and Trent (Sessions V, XIV, XXII, and XXV); Pope Benedict XII's \textit{Benedictus Deus} (DS 1000--1002); Pope St.\ Pius X's \textit{Divino afflatu} (1911); Vatican II's \textit{Gaudium et spes} (\S~18); Pope St.\ Paul VI's \textit{Laudis canticum} (1970) and \textit{Sacram unctionem infirmorum} (1972); and the \textit{Catechism of the Catholic Church} (\S\S~1005--1014, 1021--1022)---are cited strictly by canonical Denzinger-Sch\"{o}nmetzer (DS) or section locus in doctrinal paraphrase without verbatim quotation, while Fourth Lateran's \textit{Firmiter credimus} (DS 801) is checked directly in its retained Latin text.

\subsection*{Rights and Non-Authority Notice}

Offered as an independent aid for reading, reflection, and prayer, this work is not an official liturgical book, a catechism, a magisterial pronouncement, or a substitute for the sacraments and the guidance of a confessor or pastor. Project-authored exposition carries the repository license stated in the closing colophon; Scripture, official magisterial and liturgical texts, received prayers, and historical translations retain their independent public-domain or third-party rights status.
